\Def Пусть имеется последовательность операций, каждая из которых
выполняется за время $t_i$ , тогда амортизированная стоимость операции {\large $t^* = \frac{1}{n} \sum\limits_{i=1}^n t_i$}

\begin{center}
Метод монеток
\end{center}

\begin{enumerate}
    \item Допустим, что каждая операция стоит определенной количество
монет $C_i$ — учетная стоимость.
    \item Учетная стоимость не менее реальной. Тогда резерв уйдет на
покрытие остальных операций.
\end{enumerate}

Пример. Динамически расширяющийся массив
\begin{enumerate}
    \item Пусть каждый push, не требующий перевыделения памяти стоит 3 монетки. Одна монетка на запись элемента и две монетки в резерве (будем их класть на элементы массива с номерами $i$ и $i\,-\,\frac{n}{2}$ ).
    \item Инвариант. Перед каждой реаллокацией на каждом элементе есть по монете.
    \item Этими монетами оплачиваем перекопирование.
    \item То есть амортизированная сложность push равна $O(1)$.
\end{enumerate}

\pagebreak